\documentclass[10pt,a4paper]{amsart}
\usepackage[margin=19mm]{geometry}
\usepackage{amsmath,amssymb,mathtools}
\usepackage[scr=rsfs,cal=euler]{mathalfa}
\usepackage{xfrac}
\usepackage[unicode,hidelinks,bookmarks=false]{hyperref}
\newcommand{\ie}{i.e.\ }
\newcommand{\cf}{cf.\ }
\newcommand{\resp}{respectively\ }
\newcommand{\Subsection}{\subsection}
\providecommand{\nobreakdash}{\nobreak\hbox{-}}
\setlength{\emergencystretch}{2em}
\multlinegap0pt
\usepackage{amssymb}
\newtheorem{proposition}{Proposition}
\newtheorem{theorem}{Theorem}
\title{Analogue of the theta group $\Gamma_{\theta}$}
\author{Kazuhide Matsuda}
\date{Source: arXiv:2602.22471v1 (CC BY 4.0)}
\begin{document}
\maketitle

\noindent\emph{Source and licence.} Analogue of the theta group $\Gamma_{\theta}$. Version arXiv:2602.22471v1 (CC BY 4.0), distributed by arXiv under the Creative Commons Attribution 4.0 International licence, http://creativecommons.org/licenses/by/4.0/, as recorded in the arXiv licence metadata for this exact version and shown on the abstract page https://arxiv.org/abs/2602.22471v1. Full licence notice: \texttt{licenses/arxiv-2602.22471-CC-BY-4.0.txt}.\par\smallskip

\noindent\textit{Editorial scope.} Counted unit: the article's proposition prop:character-level4 (source0.tex lines 3268--3278). The article's complete original proof of it is delivered with it (source0.tex lines 3830--3886, one \texttt{proof} environment of 57 lines, which the article places away from the statement). No mathematics is written, repaired or re-derived: the statement and the proof are the article's own lines, in the article's own notation, and no retained line is substituted or re-flowed.  Retained uncounted as the source-local context the proof needs: lines 3567--3628.  Nothing was omitted from any retained span, so the package carries no omission record.  No retained line contains a citation, so the wrapper carries no bibliography and no bibliography entry.  No retained line contains a figure environment, an image-inclusion command or drawing code, so no image is delivered and the package declares no asset dependency.  Editorial formatting only, in addition: an AMS \texttt{amsart} preamble that restates the article's own declarations and macros used by the retained lines, namely the environments \texttt{proposition}, \texttt{theorem} on separate counters, together with the standard packages those lines need.  The retained environments print in this document as Proposition 1, Theorem 1 in order of appearance, whereas the article prints the same items with the numbering of its own full document.  The complete native source of the article as distributed in the arXiv e-print for this exact version is bundled unchanged as \texttt{sources/195334/source0.tex}.
\begin{theorem}
\label{thm:M:c-odd-level4}
{\it
Suppose that 
$
M=
\begin{pmatrix}
a & b \\
c & d
\end{pmatrix}
\in 
\Gamma_{\theta,4}
$ 
and 
$M\equiv 
\pm
\begin{pmatrix}
0 & -1 \\
1 & 0
\end{pmatrix}, 
%%%%%%%%%%%%%%%%%%%%%
 \pm
\begin{pmatrix}
2 & -1 \\
1 & 2
\end{pmatrix}
\bmod{4}. 
$ 
Then, 
the multiplier system 
$\nu_{G}$ of 
%%%%%%%%%%%%%%%%%%%%%%%
$
\displaystyle 
G=
\eta \left(\frac{\tau-1}{4} \right) \eta\left(\frac{\tau+1}{4} \right)
$ 
is given by 
\begin{equation*}
\nu_{G}
(M)
=
\exp
\left\{
\frac{\pi i}{6}
\left[
3(c+d-2)+\frac14(a+d)c
+
\frac14(b+c)d-cd(1+cd)
\right]
\right\}. 
\end{equation*}
In particular, we have $\nu_{G}(T)=-i,$ where 
$
T=
\begin{pmatrix}
0 & -1 \\
1 & 0
\end{pmatrix}.  
$
}
\end{theorem}

\begin{proposition}
\label{prop:character-level4}
{\it
For every $\tau\in\mathscr{H},$
set 
\begin{equation*}
G(\tau)= \eta \left(\frac{\tau-1}{4} \right) \eta\left(\frac{\tau+1}{4} \right).
\end{equation*}
The multiplier system $\nu_{G}$ is a character. 
}
\end{proposition}

\begin{proof}
By 
$
T
=
\begin{pmatrix}
0 & -1 \\
1 & 0
\end{pmatrix}, 
$ 
we have 
\begin{equation*}
MT
=
\begin{pmatrix}
a & b \\
c & d
\end{pmatrix}
%%%%%%%%%%%%%%%%%%%%%%%%%%
\begin{pmatrix}
0 & -1 \\
1 & 0
\end{pmatrix}
=
\begin{pmatrix}
b & -a \\
d & -c
\end{pmatrix}
\equiv
\pm
\begin{pmatrix}
0 & -1 \\
1 & 0
\end{pmatrix},
\pm
 \begin{pmatrix}
1 & 2 \\
2 & 1
\end{pmatrix}
\bmod{4}.
\end{equation*}
\par
Proposition \ref{prop:character-level4} and Theorem \ref{thm:M:c-odd-level4}
imply that 
\begin{align*}
\nu_{F}(MT)=&\nu_{F}(M)\nu_{F}(T)=\nu_{F}(M)\cdot(-i)  \\
=&
\exp
\left\{
\frac{\pi i}{6}
\left[
3(d-c-2)+\frac14(b-c)d+\frac14(a-d)c+cd(1-ad)
\right]
\right\},
\end{align*}
which proves the theorem. 
\end{proof}

\end{document}
